% s-313-p313-m001; original PDF 313
\Aphorism{s-313-p313-m001}{366}
Je weniger eine Philosophie von dem thatsächlich Bestehenden rechtfertigt, desto unverdächtiger ist sie. Oder genauer, um nicht dem Vorurtheil das Wort zu reden, als müsse alle Philosophie Auflehnung und Kriticismus sein: je weniger sich eine Philosophie, sei es pro oder contra, mit dem Bestehenden einlässt, desto ernsteres Gehör verdient sie. Durch vergängliche Gegner und Freunde wird man selbst um seine Ewigkeit gebracht: sehr grosse Geister, die recht gut alleinstehen konnten, sind derart in eine locale oder zeitgemässe Beschränktheit, eine Parteinahme, einen historischen Anschluss untergekrochen. Stirner gegen Feuerbach: aber was ist uns heute Feuerbach? Kant schrieb gegen Leibniz-Wolfische Curiosa, die selbst unter Theologen aus der Mode sind, verstand aber (und gab zu verstehen), er schreibe gegen Constitutionsfehler der theoretischen Vernunft. Schopenhauer hatte alle Mühe, auf keine Vorgänger und Voraussetzungen zu bauen; unter den Dingen, die er unbesehen als ewig nahm, ist aber doch etwas sehr Enges und Geschichtliches, nämlich das „metaphysische Bedürfniss“ des Menschen: in der That nur eine unberechtigte Eigenthümlichkeit des Menschen, die er sich angewöhnt hat und wieder abgewöhnen kann. Giebt es überhaupt einen unabhängigen Philosophen, der sich den Dingen, insbesondere den menschlichen Dingen gegenüber kein Für oder Wider abzwingen lässt, der nicht rechtfertigt und verdammt, der hinter das Factische von ein paar Jahrtausenden zurückdenken kann und nicht wie alle Welt das Nächste für das Wichtigste, das Auffallende für das Bedeutende hält?
\par

